% COLLECTION_METADATA: {"schema_version": 1, "kind": "reused_writeup", "independently_reviewed": false, "primary_source": "attacks/erdos/gpt_pro_5.2/410.tex", "source_paths": ["attacks/erdos/gpt_pro_5.2/410.tex"], "source_model": "gpt_pro_5.2", "source_completion": null, "source_claim": "solved"}
\section*{Erd\H{o}s Problem 410 --- collection record}

\subsection*{PROVENANCE AND STATUS}
Official problem: https://www.erdosproblems.com/410

Saved collaborative-database status: open.
Snapshot checked: 2026-09-23T17:22:24Z.
Latest status and formalization links: https://teorth.github.io/erdosproblems/

This is an attributed collection of existing writeups. The model-folder names below identify their original attribution; placement in GPT\_6\_Astra\_Ultra does not attribute their mathematical work to that model. All locally available source versions selected for this problem are reproduced below, without editing their text.

Proof claims, computations, literature statements, and completion estimates in the reused text have not been independently verified in this collection pass. A source claiming a solution does not change the saved upstream status.

Primary source (latest numbered version in the preferred model folder): \texttt{attacks/erdos/gpt\_pro\_5.2/410.tex}.

\subsection*{REUSED SOURCE 1}
Original attribution: \texttt{gpt\_pro\_5.2}.
Original file: \texttt{attacks/erdos/gpt\_pro\_5.2/410.tex}.

% BEGIN REUSED SOURCE: attacks/erdos/gpt_pro_5.2/410.tex

1) FORMAL RESTATEMENT

Let $\sigma(n)=\sum_{d\mid n} d$ be the sum-of-divisors function.
Define iterates $\sigma_1(n)=\sigma(n)$ and $\sigma_k(n)=\sigma(\sigma_{k-1}(n))$ for $k\ge 2$.

Literal question as written.
Is it true that
\[
\lim_{k\to\infty} \sigma_k(n)^{1/k}=\infty?
\]
(As written, the quantifier “for which $n$” is not explicit; the natural reading is: for every integer $n\ge 1$.)

Minimal corrected statement (standard convention).
A nontrivial version is: for every integer $n\ge 2$,
\[
\lim_{k\to\infty} \sigma_k(n)^{1/k}=\infty.
\]

2) QUICK LITERATURE/CONTEXT CHECK

The problem statement only notes that the question is discussed in Guy's collection; I do not add any external results.

3) ATTACK PLAN

Disproof-track for the literal statement:
- Check small edge cases (notably $n=1$) to see whether the limit can fail trivially.

Proof-track for the corrected statement ($n\ge 2$):
- Establish monotonic growth of $\sigma_k(n)$ and look for multiplicative lower bounds on $\sigma(n)/n$ along the orbit.

Chosen path: give a counterexample to the literal statement and briefly record basic growth facts for $n\ge 2$.

4) WORK

Counterexample to the literal statement.
Take $n=1$.
Then $\sigma(1)=1$ (the only divisor of $1$ is $1$), so the iterates satisfy
\[
\sigma_k(1)=1\quad\text{for all }k\ge 1.
\]
Therefore
\[
\sigma_k(1)^{1/k}=1^{1/k}=1\quad\text{for all }k\ge 1,
\]
and hence
\[
\lim_{k\to\infty} \sigma_k(1)^{1/k}=1\ne\infty.
\]
This disproves the literal statement “for every $n\ge 1$”.

(Additional sanity computation for $n\ge 2$.)
For example, the first few iterates for $n=2$ are
\[
2\to 3\to 4\to 7\to 8\to 15\to 24\to 60\to\cdots,
\]
which grows quickly, but this does not address the $k$th-root limit.

5) VERIFICATION

- Verified $\sigma(1)=1$ directly from the definition.
- The iteration is constant, so the limit computation is exact.

6) FINAL

**FULL SOLUTION**

**COUNTEREXAMPLE/DISPROOF**

The literal statement is false: for $n=1$, all iterates satisfy $\sigma_k(1)=1$, so $\sigma_k(1)^{1/k}=1$ for all $k$, and the limit equals $1$ rather than $\infty$.

% END REUSED SOURCE: attacks/erdos/gpt_pro_5.2/410.tex

\subsection*{COMPLETION ESTIMATE}
The primary source supplies no extractable numerical completion estimate. Its mathematical progress is not numerically assessed here.

Newly verified progress in this collection pass: $0\%$. This measures new verification work, not the amount of mathematics achieved in the reused text.
